\chapter{Grundloven og den retlige stat}

\chapterlead{Forfatningsretten handler om, hvem der må udøve offentlig magt, hvordan magten skal udøves, og hvilke grænser der beskytter borgeren. Grundloven er derfor både en organisationsplan og en beskyttelse mod vilkårlighed.}

\section{Hvorfor en grundlov?}

En almindelig lov kan normalt ændres af et nyt flertal. En grundlov har en særlig status, fordi den fastlægger de rammer, som lovgiveren selv skal arbejde inden for. Den regulerer blandt andet statens institutioner, lovgivningsproceduren, domstolene og flere grundlæggende frihedsrettigheder \citep{grundloven}.

Den danske Grundlov begynder med, at den gælder for alle dele af Danmarks Rige. Den fastlægger en indskrænket-monarkisk regeringsform og indeholder den klassiske opdeling af magten i en lovgivende, udøvende og dømmende funktion. I praksis betyder det ikke, at funktionerne er hermetisk adskilte. Regeringen forbereder lovforslag, Folketinget kontrollerer regeringen, og domstolene fortolker love. Pointen er, at ingen enkelt aktør skal kunne samle al magt uden kontrol.

\section{Magtens tredeling}

\begin{tabularx}{\textwidth}{@{}>{\RaggedRight\arraybackslash}p{2.6cm}X X@{}}
\toprule
\textbf{Funktion} & \textbf{Hovedopgave} & \textbf{Typisk kontrolspørgsmål} \\
\midrule
Lovgivende & Vedtage generelle regler og fastlægge politiske rammer. & Er loven vedtaget efter den krævede procedure og inden for Grundlovens rammer? \\
Udøvende & Gennemføre lovene og træffe konkrete afgørelser. & Har myndigheden kompetence, fulgt proceduren og anvendt skønnet sagligt? \\
Dømmende & Afgøre konkrete tvister og fortolke retten. & Er domstolen uafhængig, og har parterne fået en fair proces? \\
\bottomrule
\end{tabularx}

Opdelingen skal forstås funktionelt. En minister kan have politisk ansvar for en forvaltning, men den konkrete myndighedsafgørelse skal stadig have hjemmel og følge almindelige forvaltningsretlige krav. Domstolene kan ikke frit overtage Folketingets politiske rolle, men de skal kontrollere, om en lov eller afgørelse holder sig inden for højere regler.

\section{Legalitetsprincippet}

\term{Legalitetsprincippet} betyder i sin kerne, at offentlige indgreb over for borgere kræver et retligt grundlag. En myndighed kan ikke skabe en pligt alene ved at sige, at den synes, pligten er nyttig. Jo mere indgribende handlingen er, jo større krav stilles der typisk til hjemlens klarhed.

Princippet har flere sider:
\begin{itemize}
  \item \textbf{Kompetence:} Myndigheden skal være den rette til at træffe beslutningen.
  \item \textbf{Hjemmel:} Der skal være en regel, som bærer indgrebet.
  \item \textbf{Forudsigelighed:} Borgeren skal i rimeligt omfang kunne forstå, hvad der kræves.
  \item \textbf{Ikke-tilbagevirkning:} I strafferetten må en senere, strengere strafregel som udgangspunkt ikke anvendes på tidligere handlinger. Uden for strafferetten må spørgsmålet vurderes efter den konkrete hjemmel, ikrafttrædelsesregel og eventuelle beskyttelsesprincipper.
  \item \textbf{Kontrol:} Der skal være de kontrolmuligheder, som følger af Grundlov, lov og internationale regler; den konkrete adgang til klage eller domstolsprøvelse varierer.
\end{itemize}

Legalitetsprincippet betyder ikke, at enhver detalje skal stå direkte i en lov. Lovgivning kan delegere nærmere regulering til en minister eller myndighed. Men delegation kræver et tilstrækkeligt grundlag, og den underordnede regel må holde sig inden for rammen.

\section{Delegation og bekendtgørelser}

En lov kan for eksempel sige, at ministeren kan fastsætte nærmere regler om tekniske krav til en bygning. Det kan være praktisk, fordi tekniske standarder ændrer sig hurtigere end en lovgivningsproces. Men en bekendtgørelse kan ikke uden videre indføre en helt ny type straf, skat eller indgreb, hvis loven ikke giver hjemmel til det.

Når du læser en bekendtgørelse, skal du derfor også finde den lovbestemmelse, der bemyndiger udstedelsen. Spørg:

\begin{enumerate}
  \item Hvem har udstedt reglen?
  \item Hvilken lovbestemmelse er hjemlen?
  \item Er reglen udstedt på det rigtige område og inden for bemyndigelsen?
  \item Er der særlige krav til høring, offentliggørelse eller ikrafttræden?
\end{enumerate}

\section{Domstolenes rolle}

Domstolene afgør konkrete retstvister. De kan fortolke loven, anvende grundlovsbestemmelser og prøve, om offentlige beslutninger holder sig inden for loven. Domstolskontrol er normalt konkret: Retten tager stilling til den sag, der er forelagt, og ikke til alle tænkelige anvendelser af reglen.

Domstolene har også en institutionel begrænsning. De må ikke erstatte lovgivers politiske valg med deres egne præferencer. Men hvis et politisk valg overskrider en rettighedsbeskyttelse, en kompetenceregel eller en grundlæggende procesregel, er det netop domstolenes funktion at reagere.

\section{Grundlovens frihedsrettigheder}

Grundloven beskytter blandt andet personlig frihed, boligens ukrænkelighed, ejendomsretten, religionsfrihed, ytringsfrihed, foreningsfrihed og forsamlingsfrihed. Beskyttelsen har forskellig ordlyd og forskellig struktur. Nogle bestemmelser kræver lovhjemmel, andre indeholder særlige betingelser, og nogle skal læses sammen med menneskerettighederne.

En grundrettighed er derfor ikke altid et absolut forbud. Mange rettigheder kan begrænses, hvis betingelserne er opfyldt. Det centrale er at identificere:
\begin{itemize}
  \item hvilket beskyttet gode der berøres,
  \item om der faktisk er tale om et indgreb,
  \item om indgrebet har hjemmel,
  \item om det forfølger et legitimt formål, og
  \item om det er proportionalt og underlagt kontrol.
\end{itemize}

\section{Krise, hast og nød}

Krisebestemmelser og nødret rejser et grundlæggende spørgsmål: Kan staten handle hurtigere, når almindelige procedurer er utilstrækkelige? Svaret er ikke, at regler forsvinder, når situationen bliver vanskelig. Tværtimod bliver hjemmel, tidsbegrænsning, nødvendighed, parlamentarisk kontrol og domstolsprøvelse endnu vigtigere.

Nødret er ikke en generel tilladelse til at gøre, hvad der føles rigtigt. Den kræver en konkret vurdering af fare, nødvendighed, proportionalitet og alternativer. Hvad der var nødvendigt under en akut fare, er ikke nødvendigvis nødvendigt, når faren er aftaget.

\question{En kommune vil af sikkerhedshensyn lukke en offentlig plads efter flere urolige aftener. Hvilke spørgsmål om kompetence, hjemmel, proportionalitet, varighed og kontrol skal stilles, før du vurderer, om beslutningen er lovlig?}

\section*{Kapitelopsamling}
\begin{itemize}
  \item Grundloven fordeler magt og beskytter grundlæggende interesser.
  \item Magtens tredeling er en arbejdsdeling med gensidig kontrol, ikke tre isolerede verdener.
  \item Offentlige indgreb kræver kompetence, hjemmel, forudsigelighed og kontrol.
  \item Delegation er mulig, men en bekendtgørelse må holde sig inden for lovens ramme.
  \item Grundrettigheder kan ofte begrænses, men kun efter krav om blandt andet hjemmel og proportionalitet.
\end{itemize}
